% ============================================================
% 统一理解与生成（Unified Understanding & Generation）面试笔记
% 日期：2026/04/11
% ============================================================

\section{统一理解与生成（Unified Understanding \& Generation）}

\begin{conceptbox}
\textbf{面试速查笔记 · 2026/04/11}\\
已有基础：BAGEL MoT 架构（代码级）、Flow Matching / DDPM / DDIM 原理、BLIP3o 有所了解\\
目标：建立整个领域的技术版图，掌握流派分类、代表模型、核心设计差异
\end{conceptbox}

% -------------------- 知识地图 --------------------
\subsection{\textbf{知识地图}：你在哪里，要去哪里}

\begin{lstlisting}
[你已掌握]
  +-- BAGEL: MoT 架构 + 双支路权重 + Flow Matching 生成
  +-- Flow Matching / DDPM / DDIM 的原理与区别
  +-- BLIP3o: 知道存在，有些了解

[本次笔记覆盖]
  +-- 整个领域的"为什么要统一"动机
  +-- 四大技术流派及代表模型
  +-- 视觉表示的核心选择：离散 vs 连续
  +-- 架构分类：共享主干 vs 解耦编码
  +-- 主要模型横向对比（Chameleon / Transfusion / Janus / Show-o / BAGEL / BLIP3o）

[进阶方向]
  +-- 视频统一模型（Wan / CogVideoX-I2V）
  +-- RL 在统一模型中的应用（Flow-GRPO、UniGRPO）
  +-- Any-to-Any 生成（NExT-GPT、UnifiedIO2）
\end{lstlisting}

% ============================================================
\subsection{一、为什么需要「统一」？}
% ============================================================

\subsubsection{背景动机}

传统做法：\key{理解}和\key{生成}是两套完全独立的系统。

\begin{lstlisting}
理解侧：图片 -> ViT Encoder -> LLM -> 文本答案
生成侧：文本 -> CLIP 文本编码 -> UNet（SD/FLUX）-> 图片
\end{lstlisting}

\begin{warnbox}
\textbf{问题}：
\begin{enumerate}
  \item 两套模型维护成本高、推理成本双倍
  \item 理解能力无法增强生成（描述不好的模型也生成不好）
  \item 无法做真正的\key{交织多模态}（文字和图片穿插输入输出）
  \item 知识共享受限：图像语义理解与生成应当相辅相成
\end{enumerate}
\end{warnbox}

\subsubsection{统一的核心挑战}

\begin{intubox}
「统一」听起来直觉上很简单——把两件事塞进一个模型。但这里有一个根本矛盾：
\end{intubox}

\begin{center}
{\small
\begin{tabular}{|l|l|}
\hline
\textbf{任务} & \textbf{需要什么} \\
\hline
图像\textbf{理解} & 提取\textbf{语义}特征（SigLIP/CLIP 风格，连续高维，不可逆压缩） \\
图像\textbf{生成} & 重建\textbf{像素}细节（VAE 风格，保留纹理，可解码） \\
\hline
\end{tabular}
}
\end{center}

\begin{warnbox}
\textbf{同一个视觉表示既要理解语义又要重建像素}——这两个目标天然冲突。\\
各流派的核心分歧，本质上就是在\key{怎么解决这个矛盾}。
\end{warnbox}

% ============================================================
\subsection{二、核心维度：如何表示图像？}
% ============================================================

\subsubsection{2.1 视觉表示：离散 vs 连续}

\begin{lstlisting}
图片 --------------------------------------------------------+
      |                                                      |
      v 离散 Tokenization                          v 连续 Latent
  VQ-VAE / VQGAN                              VAE (SD 风格)
  8192 个码本索引                              连续向量（如 16x16x4）
  可以直接做 next-token 预测（AR）            需要扩散/Flow Matching 生成
  压缩率高，信息有损较多                       保留细节，但不能直接 AR
  代表：Chameleon, Janus, Show-o              代表：BAGEL, Transfusion, BLIP3o
\end{lstlisting}

\begin{conceptbox}
\textbf{直觉理解}：
\begin{itemize}
  \item \key{离散}：把图片「翻译」成一串有限词表里的单词，LLM 直接把它当文本 token 处理，整个世界只有一个序列。
  \item \key{连续}：保留图片的「模拟信号」，生成时需要扩散模型来建模连续分布，不能简单 next-token。
\end{itemize}
\end{conceptbox}

\subsubsection{2.2 架构选择：共享主干 vs 解耦编码}

\begin{lstlisting}
共享主干（Shared Backbone）         解耦编码（Decoupled Encoder）
---------------------------------   --------------------------------
理解、生成都走同一个 Transformer    理解侧用 SigLIP/CLIP（语义强）
                                    生成侧用 VQVAE/VAE（重建好）
优点：参数共享，知识互通             优点：两侧各司其职，避免冲突
缺点：两个目标相互干扰               缺点：两套编码器占显存，协调训练难
代表：Chameleon, Transfusion,       代表：Janus, Janus-Pro, EMU2
     Show-o, BAGEL(MoT)
\end{lstlisting}

% ============================================================
\subsection{三、四大技术流派}
% ============================================================

\subsubsection{流派总览}

\begin{lstlisting}
统一理解与生成
+-- 流派 A：纯 AR（离散 Token，一个序列统治一切）
|   +-- Chameleon, LlamaGen, VAR
|
+-- 流派 B：AR + 扩散混合（文本 AR，图像 Diffusion/FM）
|   +-- B1 连续 Latent：Transfusion, BAGEL, BLIP3o
|   +-- B2 离散 Masked Diffusion：Show-o
|
+-- 流派 C：解耦编码（理解 / 生成各用一套编码器）
|   +-- Janus, Janus-Pro, EMU2, SEED-X
|
+-- 流派 D：纯扩散统一（Diffusion-only）
    +-- OmniGen, InstructPix2Pix
\end{lstlisting}

% ===== 流派 A =====
\subsubsection{流派 A：纯 AR——「一切都是 Token」}

\begin{conceptbox}
\textbf{核心思想}：把图片通过 VQ-VAE/VQGAN \key{离散化}成一串 index，拼在文本 token 后面，用同一个 Transformer 做 next-token prediction。
\end{conceptbox}

\paragraph{代表：Chameleon（Meta, 2024）}

\begin{lstlisting}
[用户] "生成一只猫坐在沙发上"
  |
文本 token: [35, 821, 19, ...]
图像 token: [4521, 89, 2234, ...]（VQ-VAE 码本 index，共 1024 个）
  |
同一个 LLM，预测下一个 token（文本或图像都行）
  |
图像 token 序列 -> VQ-VAE 解码器 -> 像素
\end{lstlisting}

\textbf{训练目标}：统一的交叉熵损失，无论 token 是文字还是图像 index
$$\mathcal{L} = -\sum_{i} \log P(x_i \mid x_{<i})$$

\textbf{优点}：
\begin{itemize}
  \item 架构极简，完美的「统一」——真正只有一套模型
  \item 可以做真正的交织生成（文字中随时插图，图中随时接文字）
  \item 借助 LLM scaling law
\end{itemize}

\begin{warnbox}
\textbf{缺点}：
\begin{itemize}
  \item VQ-VAE 离散化会\warn{丢失高频细节}，生成图像质量天花板低于扩散模型
  \item 图像 token 数量多（1024+ per image），序列超长，计算量爆炸
  \item 词表从 32k 暴增到 32k + 8k（图像码本），tokenizer 需要重设计
  \item 理解和生成共用同一个 visual tokenizer，两者目标冲突
\end{itemize}
\end{warnbox}

\begin{conceptbox}
\textbf{关键技术细节}：
\begin{itemize}
  \item Chameleon 专门设计了 \key{QK-Norm} + 稳定训练策略，防止图文混合序列训练崩溃
  \item 图像 VQ-VAE 码本大小通常 8192，分辨率 $512\times512$ 对应 1024 个 token
\end{itemize}
\end{conceptbox}

\paragraph{VAR（Visual AutoRegressive, ICLR 2024 best paper）}

\note{不是统一模型，但影响了后续设计：}

\begin{lstlisting}
传统 AR：像素/patch 按光栅扫描顺序，token by token
VAR：coarse-to-fine，先生成 1x1 全局 token，再 2x2，再 4x4...
     每层预测"下一个分辨率"的所有 token（next-scale prediction）
\end{lstlisting}

\begin{summarybox}
\textbf{影响}：证明了 AR 生成可以不按光栅扫描，scale-wise 更符合图像的层次结构。
\end{summarybox}

% ===== 流派 B1 =====
\subsubsection{流派 B1：AR + 连续扩散混合}

\begin{conceptbox}
\textbf{核心思想}：文本是离散的，用 AR（next-token）；图像是连续的，用扩散/Flow Matching。两者在\key{同一个 Transformer 里}共存，但损失函数不同。
\end{conceptbox}

\begin{intubox}
这正是 \key{BAGEL} 和 \key{Transfusion} 的流派，你对 BAGEL 已经很深入了解，下面重点讲它们之间的异同。
\end{intubox}

\paragraph{代表 1：Transfusion（Meta, 2024）}

\textbf{一句话}：同一个 Transformer，文本位置算 AR 交叉熵，图像 latent 位置算扩散 MSE loss，两套 loss 联合训练。

\begin{lstlisting}
输入序列：[文本 token, ..., <image>, latent_t, ..., </image>, 文本 token, ...]
              | AR loss                | Diffusion loss
         -log P(x_i|x<i)         ||eps - eps_theta(z_t, t, context)||^2
\end{lstlisting}

\textbf{架构关键}：
\begin{itemize}
  \item 图像 latent 通过一个轻量级 \key{U-Net style 的局部卷积层}处理
  \item 文本 token 保持正常 causal attention
  \item 图像 token 之间做 \key{bidirectional attention}（扩散是并行去噪，不需要因果约束）
  \item 文本和图像 token 之间：文本单向 attend 图像，图像双向 attend 文本
\end{itemize}

\textbf{与 BAGEL 的核心区别}：

\begin{center}
{\small
\begin{tabular}{|l|p{5.5cm}|p{5.5cm}|}
\hline
 & \textbf{Transfusion} & \textbf{BAGEL} \\
\hline
图像生成范式 & DDPM 风格扩散（预测噪声 $\epsilon$） & \key{Flow Matching}（预测速度 $v$） \\
模型参数 & 单套参数（文本/图像 token 共用 Transformer 层） & \key{MoT 双套权重}（理解侧/生成侧各一套） \\
图像注意力 & 局部 U-Net 卷积 + 双向 Attention & 全局 Attention，MoT 控制参数分叉 \\
理解侧 visual encoder & 无（图像直接 patchify 进 latent） & \key{SigLIP}（独立 ViT，冻结） \\
训练时序 & 文本 AR + 图像扩散联合训练 & 同上，但有 Flow Matching 时间步 \\
\hline
\end{tabular}
}
\end{center}

\paragraph{代表 2：BAGEL（字节, 2025）}

\begin{intubox}
关键复习：BAGEL 的核心创新是 \key{MoT（Mixture of Transformers）}——同一层里，文本 token 走理解支路（无后缀），图像 latent 走生成支路（\code{moe\_gen}），一次 Attention 完成跨模态交互。
\end{intubox}

\begin{lstlisting}
Transfusion：单套参数，试图让同一套权重同时理解和生成
    | 发现两个目标互相干扰
BAGEL：物理分开两套权重（MoT），但共享一次 Attention
    -> 理解侧专心学理解，生成侧专心学生成
    -> 跨模态语义通过 Attention 自然流通
\end{lstlisting}

\textbf{BAGEL 的完整技术栈}：
\begin{lstlisting}
图像理解侧：图片 -> SigLIP ViT -> patch embed -> LLM (理解支路)
图像生成侧：VAE latent (带噪) -> MLP projector -> LLM (生成支路)
文本：tokenizer -> embedding -> LLM (理解支路)
生成输出：LLM (生成支路) -> MLP projector -> VAE latent -> VAE 解码 -> 图片
训练目标：理解侧 AR 交叉熵 + 生成侧 Flow Matching MSE
\end{lstlisting}

\paragraph{代表 3：BLIP3o（Salesforce, 2024）}

\textbf{定位}：偏向图像生成质量的统一模型，继承 BLIP 系列的理解强项，加入 Flow Matching 生成。

\textbf{核心架构}：
\begin{lstlisting}
理解侧：图片 -> SigLIP / EVA-CLIP -> Q-Former or MLP Projector -> LLM
生成侧：文本 -> LLM -> 隐层特征 -> Flow Matching Decoder -> 图片
\end{lstlisting}

\textbf{关键技术}：
\begin{itemize}
  \item \key{Diffusion Head / Flow Matching Head}：独立的扩散解码器，接收 LLM 最后一层的条件向量
  \item 与 BAGEL 不同：图像 latent \textbf{不进入} LLM 序列里做联合建模，而是 LLM 先编码条件，扩散头单独生成
  \item 更像：「LLM as Text Encoder + 独立的 Diffusion Decoder」
\end{itemize}

\begin{lstlisting}
BAGEL：[文本 token, 图像 latent token] 在同一个 LLM 序列里 -> 更深的融合
BLIP3o：LLM 处理文本/理解 -> 输出条件向量 -> 独立扩散头生成 -> 更模块化
\end{lstlisting}

\begin{summarybox}
\textbf{面试要点}：BLIP3o 代表了一类「轻度统一」的思路——LLM 主体不动，通过插拔生成头来加入生成能力，实现上更简单，但理解与生成的交互深度不如 BAGEL/Transfusion。BLIP3o 的图像 VQA / 描述能力强，是它的相对优势。
\end{summarybox}

% ===== 流派 B2 =====
\subsubsection{流派 B2：AR + 离散掩码扩散}

\paragraph{代表：Show-o（NUS, 2024）}

\textbf{一句话}：文本用 AR（自回归），图像用 \key{Masked Diffusion Language Model（MDLM）}——同一个 Transformer，两种 loss。

\textbf{关键设计}：
\begin{lstlisting}
文本 token：causal autoregressive，从左到右预测
图像 token (VQ-VAE 离散)：masked diffusion，随机 mask 掉图像 token，
                          预测被 mask 的位置
                          -> 类似 BERT 的 MLM，但有时间步 t 控制 mask 比例
\end{lstlisting}

\begin{conceptbox}
\textbf{与 Chameleon 的区别}：
\begin{itemize}
  \item Chameleon：图像也是 AR（从左到右）
  \item Show-o：图像是\key{双向}的 masked diffusion（所有位置并行预测，不是序列生成）
\end{itemize}
\textbf{与 Transfusion/BAGEL 的区别}：
\begin{itemize}
  \item Show-o 的图像 token 仍是\key{离散}的（VQ-VAE index）
  \item Transfusion/BAGEL 用\key{连续} latent + 扩散
\end{itemize}
\end{conceptbox}

\begin{intubox}
\textbf{直觉}：Show-o 是「离散 token 界的妥协方案」——不想像 Chameleon 那样让图像也做 AR（太慢，质量差），但又不想引入连续扩散（太复杂），于是用离散掩码扩散——速度快、并行、还能用码本的整数表示。
\end{intubox}

% ===== 流派 C =====
\subsubsection{流派 C：解耦视觉编码}

\begin{conceptbox}
\textbf{核心动机}：与其让一个 visual tokenizer 同时服务理解和生成（根本矛盾），不如\key{两件事各用各的工具}。
\end{conceptbox}

\paragraph{代表：Janus（DeepSeek, 2024）\& Janus-Pro}

\begin{lstlisting}
输入图片（理解时）-> SigLIP ViT（语义强，压缩激进）-> LLM
输出图片（生成时）-> LLM -> VQ-VAE decoder（细节保留好）-> 图片
              ^
          同一个 LLM 主干
\end{lstlisting}

\textbf{关键设计}：
\begin{itemize}
  \item 理解时走 \key{SigLIP}（高语义，$4\times4$ 下采样，每张图 $\sim$256 token）
  \item 生成时走 \key{VQVAE}（高保真，用于重建，1024 token per image）
  \item 两套编码器共享同一个 LLM Backbone
  \item 没有像 BAGEL 那样做双套权重——LLM 参数共享，但输入来源不同
\end{itemize}

\textbf{Janus-Pro 改进}：
\begin{itemize}
  \item 更大的图像 tokenizer（码本 16384 vs 8192）
  \item 更多样的训练数据
  \item 训练策略三阶段：纯理解 $\to$ 纯生成 $\to$ 联合微调
\end{itemize}

\textbf{与 BAGEL 的本质区别}：

\begin{center}
{\small
\begin{tabular}{|l|p{5.5cm}|p{5.5cm}|}
\hline
 & \textbf{Janus} & \textbf{BAGEL} \\
\hline
理解侧编码 & SigLIP（进 LLM 的 Q/K/V 统一参数） & SigLIP（进 LLM 的\key{理解支路}专属参数） \\
生成侧编码 & VQVAE 离散 token，AR 预测 & VAE 连续 latent，Flow Matching \\
LLM 内部 & \key{共享}一套参数（理解/生成 token 共用同一套 Q/K/V） & \key{双套}参数（MoT，理解/生成各一套） \\
生成质量 & 受限于 VQ-VAE 离散化 & 连续 latent + FM，质量更高 \\
理解干扰 & 理解/生成共用 LLM 参数，有干扰 & MoT 分支，减少干扰 \\
\hline
\end{tabular}
}
\end{center}

\begin{intubox}
\textbf{直觉}：Janus 的解耦只做了「一半」——输入侧解耦了（不同编码器），但 LLM 内部仍是共享参数，理解和生成还是会抢参数空间。BAGEL 的 MoT 把解耦做到了 LLM 内部每一层。
\end{intubox}

% ===== 流派 D =====
\subsubsection{流派 D：纯扩散统一}

\paragraph{代表：OmniGen（2024）}

\textbf{思路}：不用 AR 语言模型，直接用一个扩散模型处理\textbf{所有输入}（包括文字 condition 和图像 condition），输出图像。

\begin{lstlisting}
任意输入（文字 + 图片 + 指令）-> Encoder -> 条件向量
         |
 DiT (Diffusion Transformer) -> 去噪 -> 输出图像
\end{lstlisting}

\textbf{适用场景}：图像编辑、风格迁移、条件生成——一个模型搞定。

\begin{warnbox}
\textbf{局限}：输出只有图像，不能输出文字，所以「理解」能力弱（没有文字输出接口）。
\end{warnbox}

% ============================================================
\subsection{四、主流模型横向对比表}
% ============================================================

\begin{table}[h]
\centering
\small
\begin{tabular}{|p{1.8cm}|p{1.5cm}|p{2cm}|p{2cm}|p{2cm}|p{2cm}|p{3cm}|}
\hline
\textbf{模型} & \textbf{机构} & \textbf{图像表示} & \textbf{生成范式} & \textbf{LLM 参数} & \textbf{理解编码器} & \textbf{架构特点} \\
\hline
\textbf{Chameleon} & Meta & 离散 VQ & AR next-token & 共享 & 无独立编码器（VQ-VAE 做理解） & 极简统一，质量受限 \\
\hline
\textbf{Transfusion} & Meta & 连续 latent & DDPM 扩散 & 共享（单套） & 无独立 ViT & 文本 AR + 图像扩散，联合 loss \\
\hline
\textbf{Show-o} & NUS & 离散 VQ & Masked Diffusion & 共享 & 无（VQ token 直接输入） & AR 文本 + 掩码扩散图像 \\
\hline
\textbf{Janus} & DeepSeek & 离散 VQ & AR next-token & 共享 & SigLIP（理解侧专用） & 解耦输入编码器 \\
\hline
\textbf{Janus-Pro} & DeepSeek & 离散 VQ & AR next-token & 共享 & SigLIP & Janus 升级版，更大码本 \\
\hline
\textbf{BAGEL} & 字节 & 连续 VAE & \textbf{Flow Matching} & \textbf{MoT 双套} & SigLIP（冻结 ViT） & MoT，文本 AR + latent FM \\
\hline
\textbf{BLIP3o} & Salesforce & 连续 VAE & Flow Matching & 共享（外挂 FM 头） & SigLIP / EVA-CLIP & LLM 输出条件 $\to$ 独立扩散头 \\
\hline
\textbf{EMU2} & BAAI & 连续 & AR + 扩散 & 共享 & EVA-CLIP & 多模态交织生成 \\
\hline
\textbf{OmniGen} & — & 连续 & 纯 DiT 扩散 & 无 LLM & — & 全扩散，无文字输出 \\
\hline
\end{tabular}
\end{table}

% ============================================================
\subsection{五、视觉 Tokenizer 深挖（面试常考）}
% ============================================================

\subsubsection{5.1 VQ-VAE vs VAE：两套逻辑}

\begin{lstlisting}
VQ-VAE（离散）：
图片 -> Encoder -> 连续向量 -> 最近邻码本查找 -> 离散 index
               <------------ 训练反传（straight-through estimator）
离散 index -> 码本 embed -> Decoder -> 重建图片

VAE（连续）：
图片 -> Encoder -> mu, sigma -> 采样 z = mu + sigma * eps
z -> Decoder -> 重建图片
训练：重建 loss + KL 散度
\end{lstlisting}

\begin{conceptbox}
\textbf{关键区别}：
\begin{itemize}
  \item VQ-VAE 输出整数 index（可以直接被 LLM 的 vocab 处理）
  \item VAE 输出连续向量（不能 next-token prediction，必须用扩散）
\end{itemize}
\end{conceptbox}

\subsubsection{5.2 SigLIP 为什么只用于理解不用于生成？}

SigLIP（Sigmoid Loss Image-Language Pretraining）是 OpenCLIP 的改进版：
\begin{itemize}
  \item 训练目标：让图文对的 sigmoid 相似度高，非对应的低
  \item 输出：高度\key{语义化}的 embedding，丢掉了像素级细节
  \item 压缩激进：一张 $224\times224$ 图片 $\to$ 196 个 patch token（$14\times14$）
\end{itemize}

\textbf{对比 VAE（用于生成）}：
\begin{itemize}
  \item VAE 训练目标是\key{重建像素}，保留高频细节
  \item 一张图 $\to$ $16\times16 = 256$ 个 latent，每个是 4/8/16 维连续向量
\end{itemize}

\begin{lstlisting}
SigLIP -> 告诉你"图里有猫" -> 适合理解（VQA、Caption）
VAE    -> 告诉你"猫的每根毛的颜色" -> 适合生成（重建）
\end{lstlisting}

\begin{summarybox}
BAGEL、Janus 都用 SigLIP 做理解侧，VAE/VQVAE 做生成侧，正是认识到了这个根本区别。
\end{summarybox}

% ============================================================
\subsection{六、训练策略：统一模型怎么练？}
% ============================================================

\subsubsection{6.1 多阶段训练（几乎所有模型都用）}

\begin{lstlisting}
阶段 1（预训练基础）：大规模图文对，分别练理解和生成，避免早期干扰
  |
阶段 2（联合训练）：理解任务 + 生成任务同时训练，loss 加权求和
  |
阶段 3（指令微调 SFT）：高质量指令数据，对齐人类意图
  |
阶段 4（可选 RL 微调）：GRPO/DPO，进一步提升生成质量（如 BAGEL 的 Flow-GRPO）
\end{lstlisting}

\subsubsection{6.2 Loss 平衡是核心挑战}

\begin{warnbox}
\textbf{问题}：
\begin{itemize}
  \item 理解任务的 loss（交叉熵）和生成任务的 loss（Flow Matching MSE）量级不同
  \item 如果直接加权，某一侧会 dominate，另一侧 loss 退化
\end{itemize}
\end{warnbox}

\textbf{常见做法}：
\begin{itemize}
  \item \key{Loss 归一化}：分别对两侧 loss 归一化到同一量级
  \item \key{Gradient surgery}：梯度冲突时投影，避免两侧任务互相破坏
  \item \key{Batch ratio}：控制理解样本 vs 生成样本的比例
  \item \key{训练阶段分离}：早期只练一侧，稳定后才联合
\end{itemize}

\subsubsection{6.3 Classifier-Free Guidance（CFG）在统一模型中}

生成图片时通常用 CFG 提升质量：
$$\hat{\epsilon} = \epsilon_{\text{uncond}} + w \cdot (\epsilon_{\text{cond}} - \epsilon_{\text{uncond}})$$

\begin{conceptbox}
\textbf{统一模型的 CFG 实现}：
\begin{itemize}
  \item \textbf{条件}：文本描述（通过 LLM 编码）
  \item \textbf{无条件}：用空字符串/null token 作为条件
  \item BAGEL 的 \code{InterleaveInferencer} 中维护两份 context（有条件 + 无条件），CFG 在 Flow Matching 去噪时应用
\end{itemize}
\end{conceptbox}

% ============================================================
\subsection{七、面试高频问题 \& 答题框架}
% ============================================================

\subsubsection{Q1：统一理解与生成的核心挑战是什么？}

\begin{summarybox}
\textbf{答题框架}：核心矛盾是\key{视觉表示的双重需求}：理解需要语义压缩（SigLIP 风格），生成需要像素重建（VAE 风格）。各流派的核心分歧就是怎么解决这个矛盾：
\begin{itemize}
  \item Chameleon/Janus 选\key{两套编码器}（VQ 理解 + VQ 生成或 CLIP + VQ）
  \item BAGEL 选\key{两套编码器 + MoT 双套 LLM 权重}（把分离推进到 LLM 内部）
  \item Transfusion 选\key{一套权重}但接受两侧相互干扰，用联合 loss 平衡
  \item BLIP3o 选\key{外挂生成头}，最小侵入 LLM 主体
\end{itemize}
\end{summarybox}

\subsubsection{Q2：BAGEL 和 Janus 都用了 SigLIP，有什么本质区别？}

\begin{summarybox}
\textbf{答}：
\begin{itemize}
  \item \textbf{Janus}：SigLIP patch 进入 LLM 后，理解和生成 token 共用\key{同一套} Q/K/V/MLP，LLM 内部没有分离
  \item \textbf{BAGEL}：SigLIP patch 进入 LLM 的\key{理解支路}（无后缀参数），图像 VAE latent 进入\key{生成支路}（\code{moe\_gen} 后缀参数），同一层内两套权重，但一次 Attention 互通
\end{itemize}
本质区别：BAGEL 的分离是\key{参数级}的（每层都有两套），Janus 的分离只是\key{编码器级}的（LLM 以前分，LLM 里共用）。
\end{summarybox}

\subsubsection{Q3：为什么 BAGEL 用 Flow Matching 而不是 DDPM？}

\begin{summarybox}
\textbf{答}：
\begin{itemize}
  \item Flow Matching 路径是线性插值（直线），DDPM 是弯曲随机游走
  \item 训练更稳定：FM 的目标 $v = x_1 - x_0$ 是常数方向，不同时间步的 loss 量级均匀
  \item 采样更快：直线路径让 ODE 积分步数更少（与 FLUX/SD3 同一技术路线）
  \item BAGEL 联合 LLM 训练时，训练稳定性很重要，FM 的均匀 loss 有助于多任务平衡
\end{itemize}
\end{summarybox}

\subsubsection{Q4：离散 Token（VQ-VAE）和连续 Latent（VAE）各有什么优劣？}

\begin{center}
{\small
\begin{tabular}{|l|p{5cm}|p{5cm}|}
\hline
 & \textbf{离散 VQ-VAE} & \textbf{连续 VAE} \\
\hline
与 LLM 的兼容性 & 天然兼容（整数 index） & 需要额外扩散模型 \\
生成质量上限 & 受码本大小限制，有量化误差 & 更高，无量化损失 \\
生成速度 & AR 逐 token 慢（1024+ 步） & 扩散并行，通常更快 \\
训练复杂度 & 统一 loss（交叉熵），简单 & 多种 loss，需平衡 \\
代表模型 & Chameleon, Janus, Show-o & BAGEL, Transfusion, BLIP3o \\
\hline
\end{tabular}
}
\end{center}

\begin{summarybox}
\textbf{趋势}：工业界高质量生成正在转向\key{连续 latent + Flow Matching}（与 FLUX/SD3 同轨）。
\end{summarybox}

\subsubsection{Q5：统一模型如何处理图像理解中的「多图输入」问题？}

\begin{summarybox}
\textbf{答}：以 BAGEL 为例：
\begin{enumerate}
  \item 每张图片通过 SigLIP ViT 提取 patch features（约 256--1024 个 token/图）
  \item 通过 MLP projector 映射到 LLM 的隐藏维度
  \item 拼入序列，做 \code{forward\_cache\_update\_vit}（更新 KV cache）
  \item 多张图片顺序缓存，形成长上下文，LLM 统一理解
\end{enumerate}
\warn{挑战}：多图会导致序列极长（4 张图 $\times$ 1024 token = 4096 仅用于图像），需要 flash attention 或 sequence packing。
\end{summarybox}

\subsubsection{Q6：Chameleon 训练会不稳定，为什么？怎么解决？}

\begin{summarybox}
\textbf{答}：
\begin{itemize}
  \item \textbf{原因}：文字 token 和图像 token 的 embedding 分布差异大，联合训练时梯度冲突，导致 loss spike 甚至训练崩溃
  \item \textbf{Chameleon 的解决方案}：\key{QK-Norm}（对 Attention 的 Query 和 Key 做 RMSNorm），控制内积量级，稳定 softmax
  \item 此外：\key{z-loss}（对 logits 做正则），防止某些 token 概率接近 1
\end{itemize}
\end{summarybox}

\subsubsection{Q7：什么是 MoT（Mixture of Transformers），它解决了什么问题？}

\begin{summarybox}
\textbf{答}：MoT 是 BAGEL 的核心架构创新。传统统一模型让理解和生成 token 共用同一套参数，导致两个目标互相干扰。MoT 在每个 Transformer 层内物理分为两套参数（理解支路无后缀，生成支路 \code{moe\_gen}），但\key{只做一次 Attention}——两侧 token 在注意力中仍然互相可见，语义交流没有割断。

\textbf{三要素}：一条序列（文本 + 图像 latent 打包）、两套权重（按 token 类型切换）、一次 Attention（全局互通）。
\end{summarybox}

% ============================================================
\subsection{八、知识体系图：从 BAGEL 出发看整个领域}
% ============================================================

\begin{lstlisting}
                    统一理解与生成
                         |
          +--------------+--------------+
          v              v              v
    视觉表示         LLM 架构        生成范式
    离散 vs 连续     共享 vs 解耦    AR vs 扩散 vs 混合
\end{lstlisting}

\begin{summarybox}
你的位置：BAGEL（连续 + 解耦 MoT + 混合 Flow Matching）是最前沿的设计之一。
\end{summarybox}

% ============================================================
\subsection{九、推荐论文}
% ============================================================

\begin{center}
{\small
\begin{tabular}{|l|l|p{9cm}|}
\hline
\textbf{论文} & \textbf{时间} & \textbf{核心贡献} \\
\hline
Chameleon & 2024.05 & Meta 首个真正统一的 AR 多模态模型 \\
Transfusion & 2024.08 & 文本 AR + 图像扩散，联合训练范式 \\
Janus / Janus-Pro & 2024.10 / 2025.01 & 解耦视觉编码，DeepSeek \\
Show-o & 2024.09 & AR + 离散掩码扩散混合 \\
BAGEL & 2025 & MoT + Flow Matching \\
BLIP3o & 2024 & LLM + 独立扩散头，Salesforce \\
VAR & 2024 & Next-scale AR，ICLR best paper \\
\hline
\end{tabular}
}
\end{center}

% ============================================================
\subsection{补充高频面试题（来源：知乎 / 牛客 / CSDN）}
% ============================================================

% ------- Category 1: 统一模型动机与核心挑战 -------
\subsubsection{一、统一模型动机与核心挑战}

\begin{conceptbox}{Q1-1：为什么"理解"与"生成"难以统一？三大核心矛盾是什么？}
\key{三大矛盾}：
\begin{enumerate}
  \item \textbf{表示空间}：理解需要\key{语义压缩}的高层特征（SigLIP/CLIP），生成需要\key{像素级细节}的连续隐变量（VAE）——同一编码器无法同时满足。
  \item \textbf{注意力模式}：文本生成用\key{causal mask}（前向自回归），图像理解/扩散去噪用\key{bidirectional attention}（全局上下文）——强行混用会破坏自回归的时序因果性。
  \item \textbf{训练目标}：理解优化\key{交叉熵}（离散分类），生成优化\key{MSE / Flow Matching}（连续回归）——损失量纲和梯度尺度差异悬殊，联合训练极易失衡。
\end{enumerate}
\end{conceptbox}

\begin{conceptbox}{Q1-2：从信息论角度，为什么理解编码器不能直接用于生成？}
理解编码器（如 SigLIP）通过\key{对比学习}训练：
\[
\mathcal{L}_{\text{CLIP}} = -\log \frac{\exp(\text{sim}(v_i, t_i)/\tau)}{\sum_j \exp(\text{sim}(v_i, t_j)/\tau)}
\]
对比损失只保留\key{跨模态判别信息}，丢弃了大量\warn{像素级低频细节}（颜色、纹理、空间布局）。

而生成任务的解码器需要从隐变量\textbf{重建}完整像素——若隐变量缺失像素细节，解码器无从恢复。

\good{VAE} 通过重建损失 $\mathcal{L}_{\text{recon}} = \|x - \hat{x}\|^2$ 强制保留像素信息，因此适合生成端。
\end{conceptbox}

\begin{conceptbox}{Q1-3：统一模型的三大技术路线及其核心取舍}
\begin{center}
\small
\begin{tabular}{|l|l|l|l|}
\hline
\textbf{路线} & \textbf{代表模型} & \textbf{优势} & \textbf{劣势} \\
\hline
纯 AR（离散 token） & Chameleon, VAR & 架构统一，推理简单 & 图像质量受 VQ 瓶颈限制 \\
AR + 扩散混合 & Transfusion, BAGEL & 文本 AR 质量 + 图像扩散质量 & 训练/推理逻辑复杂 \\
解耦编码器 & Janus, BLIP3o & 各模块专注各自目标 & 参数冗余，联动性弱 \\
\hline
\end{tabular}
\end{center}
\end{conceptbox}

\begin{conceptbox}{Q1-4：统一模型相比"理解模型 + 生成模型"拼接的优势？}
\begin{itemize}
  \item \good{共享世界知识}：LLM 中积累的常识、逻辑推理能力可直接增强图像生成的语义一致性。
  \item \good{减少系统复杂度}：单模型推理，无需调度两套系统，延迟更低。
  \item \good{多模态对齐}：理解与生成共享 token 空间，文图互为上下文，有助于精细的指令跟随生成。
  \item \warn{代价}：训练数据需求量大，多任务权衡（\key{tug-of-war}）问题——理解性能可能因生成目标而下降。
\end{itemize}
\end{conceptbox}

% ------- Category 2: 视觉表示选择 -------
\subsubsection{二、视觉表示选择：离散 Token vs 连续隐变量}

\begin{conceptbox}{Q2-1：VQ-VAE vs VAE 全面对比}
\begin{center}
\small
\begin{tabular}{|l|l|l|}
\hline
\textbf{维度} & \textbf{VQ-VAE（离散）} & \textbf{VAE（连续）} \\
\hline
隐空间 & 离散 codebook 索引 & 连续高斯分布 $\mathcal{N}(\mu, \sigma^2)$ \\
梯度 & STE 直通估计（有偏） & 重参数化（无偏） \\
与 LLM 融合 & 天然契合（token 序列） & 需额外离散化或扩散头 \\
图像质量上限 & 受 codebook 大小限制 & 理论无损，质量更高 \\
推理速度 & AR 逐 token，较慢 & 扩散多步，可并行化 \\
代表模型 & Chameleon, Show-o, Janus（生成端） & Transfusion, BAGEL, BLIP3o \\
\hline
\end{tabular}
\end{center}
\end{conceptbox}

\begin{conceptbox}{Q2-2：VQ-VAE 的完整损失函数是什么？STE 的作用？}
\[
\mathcal{L}_{\text{VQ-VAE}} = \underbrace{\|x - \hat{x}\|^2}_{\text{重建损失}} + \underbrace{\|\text{sg}[z_e] - e\|^2}_{\text{codebook 损失}} + \underbrace{\beta \|z_e - \text{sg}[e]\|^2}_{\text{commitment 损失}}
\]
其中 $z_e = E(x)$，$e = \arg\min_k \|z_e - e_k\|$，$\text{sg}[\cdot]$ 为 stop-gradient 算子。

\key{STE（Straight-Through Estimator）}：前向传播使用量化值 $e$，\textbf{反向传播直接将梯度复制给} $z_e$（即 $\frac{\partial e}{\partial z_e} \approx 1$）。这是一个\warn{有偏估计}，但实践中有效，是离散化可微训练的关键技巧。
\end{conceptbox}

\begin{conceptbox}{Q2-3：Codebook Collapse 是什么？如何缓解？}
\warn{Codebook Collapse}：训练中大量 codebook entry 从未被使用，仅少数 entry 被频繁查询，导致表示能力退化。

\textbf{根本原因}：编码器输出分布集中，远离未激活 entry；$\arg\min$ 操作使梯度不流向冷门 entry。

\textbf{缓解方法}：
\begin{itemize}
  \item \good{EMA 更新 codebook}（而非梯度）：$e_k \leftarrow \alpha e_k + (1-\alpha) \bar{z}_e^{(k)}$
  \item \good{Random restart}：将使用率过低的 entry 重置为近期编码器输出
  \item \good{LFQ（Lookup-Free Quantization，MAGVIT-v2）}：将特征向量每个维度\textbf{独立二值化}（$\pm 1$），codebook 大小指数级扩展至 $2^d$，无需显式 lookup，彻底消除 collapse
\end{itemize}
\end{conceptbox}

\begin{conceptbox}{Q2-4：为什么 Janus 的理解端用 SigLIP 而生成端用 VQVAE？}
\begin{itemize}
  \item \textbf{SigLIP for 理解}：Sigmoid 形式的图文对比学习，输出\key{高层语义特征}，与文本表示空间对齐，适合 VQA / 描述 / 推理任务。但它\warn{丢失像素细节}，无法重建图像。
  \item \textbf{VQVAE for 生成}：离散 token 序列可直接被 AR Transformer 建模（next-token prediction），训练目标统一为 CE 损失，无需额外的扩散推理模块。
\end{itemize}
\key{核心洞见}：理解和生成对"好的视觉表示"的需求\textbf{完全相反}——解耦是目前最稳健的工程选择。
\end{conceptbox}

% ------- Category 3: 纯 AR 流派 -------
\subsubsection{三、纯 AR 流派（Chameleon / VAR）}

\begin{conceptbox}{Q3-1：Chameleon 训练不稳定的根本原因及 QK-Norm 的解决方案}
\warn{问题根源}：混合模态训练中，图像 token（高密度视觉信息）与文本 token 的特征分布\textbf{量纲差异极大}，导致 Attention 的 $QK^\top$ 内积出现数值爆炸（softmax 饱和 $\to$ 梯度消失/爆炸）。

\key{QK-Norm 方案}（来自 Chameleon 论文）：
\[
\text{Attention}(Q, K, V) = \text{softmax}\!\left(\frac{\text{norm}(Q) \cdot \text{norm}(K)^\top}{\tau}\right) V
\]
对 $Q$ 和 $K$ 分别做 LayerNorm，将内积值约束在 $[-1, 1]$ 范围内，\good{彻底消除数值爆炸}，训练稳定性显著提升。

\note{这一技巧已被后续多模态模型（Transfusion, BAGEL 等）广泛采用。}
\end{conceptbox}

\begin{conceptbox}{Q3-2：VAR（Visual AutoRegressive）的 next-scale 预测与标准 next-token 有何不同？}
\begin{itemize}
  \item \textbf{标准 AR（Chameleon）}：将图像展平为 1D token 序列，按\key{光栅扫描顺序}逐 token 预测，序列长度 $\sim 256$，推理慢。
  \item \textbf{VAR}：将图像表示为\key{多尺度 token map}（$1^2 \to 2^2 \to \cdots \to 16^2$），每步预测\textbf{下一个分辨率}的完整 token map，利用低分辨率全局结构指导高分辨率细节。
\end{itemize}
\good{优势}：FID=1.73（超越 DiT），推理步数 = 尺度数（约 10 步），速度快 $\sim 20\times$ vs 同分辨率 AR。

\warn{局限}：仍依赖 VQ 离散化，高频细节受 codebook 质量限制。
\end{conceptbox}

\begin{conceptbox}{Q3-3：纯 AR 图像生成相比扩散模型的质量瓶颈}
\begin{center}
\small
\begin{tabular}{|l|l|l|}
\hline
\textbf{方面} & \textbf{纯 AR} & \textbf{扩散模型} \\
\hline
隐空间 & 离散 VQ，有量化误差 & 连续 VAE，信息保留更完整 \\
生成顺序 & 强制因果依赖，全局一致性差 & 全局并行去噪，一致性好 \\
CFG 强度 & 有限（token 分类 logit 调整） & 连续空间，强度可任意调节 \\
FID & VAR 1.73（SOTA AR） & DiT / SD 系列 $<1.5$ 可达 \\
\hline
\end{tabular}
\end{center}
\end{conceptbox}

% ------- Category 4: AR + 扩散混合 -------
\subsubsection{四、AR + 扩散混合流派（Transfusion / BAGEL / BLIP3o）}

\begin{conceptbox}{Q4-1：Transfusion 的双向注意力设计——为什么图像 patch 用 bidirectional attention？}
Transfusion 在同一 Transformer 中混合两种注意力掩码：
\begin{itemize}
  \item \textbf{文本 token}：\key{causal mask}，维持自回归性质，支持 next-token 预测。
  \item \textbf{图像 patch（扩散）}：\key{bidirectional（full）attention}，图像内所有 patch 互相可见。
\end{itemize}
\key{原因}：扩散去噪是\textbf{全局降噪}过程——每个 patch 的去噪方向取决于整张图的上下文。若使用 causal mask，patch 只能看到左侧 patch，全局结构（如对称布局、颜色分布）无法被捕捉，去噪质量严重下降。

\note{文本与图像 patch 之间：文本可 attend 到所有先前 token，图像 patch 可 attend 到所有文本 token + 同图所有 patch。}
\end{conceptbox}

\begin{conceptbox}{Q4-2：BAGEL 的 MoT（Mixture of Transformers）原理}
\key{MoT 核心设计}：
\begin{itemize}
  \item 所有 token（文本 / 图像）\textbf{共享同一个 Attention 层}（跨模态交互）
  \item FFN 层\textbf{按 token 类型路由}：文本 token $\to$ FFN$_{\text{text}}$；图像 token $\to$ FFN$_{\text{image}}$
\end{itemize}

\textbf{对比 MoE}：MoE 用\warn{学习的路由器}（softmax + top-k）动态分配 token，MoT 用\good{确定性类型标签}（完全无额外参数开销）。

\textbf{优势}：
\begin{itemize}
  \item 参数\textbf{专用化}（FFN 各司其职），无模态干扰
  \item Attention 层跨模态共享，保留文图联合理解能力
  \item 相比独立双模型，参数量减少约 40\%，无负载均衡 loss
\end{itemize}

\warn{注意}：MoT 中没有路由概率，无法做 soft mixing，适合模态\textbf{明确可分}的场景。
\end{conceptbox}

\begin{conceptbox}{Q4-3：BLIP3o 的架构与 Transfusion/BAGEL 的关键区别}
\begin{center}
\small
\begin{tabular}{|l|l|l|l|}
\hline
& \textbf{BLIP3o} & \textbf{Transfusion} & \textbf{BAGEL} \\
\hline
LLM 与扩散关系 & \textbf{完全解耦}，独立扩散头 & 同一 Transformer，混合 loss & MoT 共享 Attn \\
图像隐变量来源 & LLM 输出 CLIP 特征 & VAE 连续 latent & VAE 连续 latent \\
扩散目标 & 在 CLIP 特征空间扩散 & 在像素空间扩散 & Flow Matching（RF）\\
优势 & LLM 可换，扩散头可换 & 端到端联合优化 & 模态专用 FFN \\
劣势 & 两段训练，CLIP 特征损失像素细节 & 架构耦合，扩展性差 & 路由硬编码，不够灵活 \\
\hline
\end{tabular}
\end{center}
\end{conceptbox}

% ------- Category 5: 解耦视觉编码 -------
\subsubsection{五、解耦视觉编码（Janus / Janus-Pro）}

\begin{conceptbox}{Q5-1：Janus 架构的完整数据流}
\textbf{理解路径}：输入图像 $\xrightarrow{\text{SigLIP}}$ 高维语义特征 $\xrightarrow{\text{Adapter}}$ LLM token $\to$ 文本输出

\textbf{生成路径}：文本 prompt $\to$ LLM $\to$ VQVAE 离散 token 序列 $\xrightarrow{\text{VQVAE decoder}}$ 重建图像

两条路径\key{共享 LLM backbone}，但视觉编码器\textbf{完全独立}，互不干扰。

\textbf{训练}：理解阶段用图文对 SFT（冻结 SigLIP），生成阶段用 AR next-token CE（冻结 VQVAE 编码器和解码器）。
\end{conceptbox}

\begin{conceptbox}{Q5-2：Janus vs BAGEL 对比}
\begin{center}
\small
\begin{tabular}{|l|l|l|}
\hline
\textbf{维度} & \textbf{Janus} & \textbf{BAGEL} \\
\hline
图像生成方式 & VQVAE AR（离散） & Flow Matching（连续） \\
图像质量 & 受 VQ 量化限制 & 更高（连续隐空间） \\
理解编码器 & SigLIP（冻结） & SigLIP（冻结） \\
生成编码器 & VQVAE（冻结） & VAE（冻结） \\
架构统一性 & 低（两个独立编码器） & 高（MoT 统一 backbone） \\
推理复杂度 & 低（纯 AR） & 高（AR + 多步 FM） \\
\hline
\end{tabular}
\end{center}
\end{conceptbox}

\begin{conceptbox}{Q5-3：Janus-Pro 在 Janus 基础上的改进点}
\begin{itemize}
  \item \good{更大的生成 VQ 码本}：从 16384 扩展至更大 codebook，减少量化误差，图像细节更丰富。
  \item \good{训练数据扩展}：加入更多高质量图文生成数据，GenEval 分数从 0.61 提升至 0.80+。
  \item \good{解码器升级}：VQVAE 解码器换用更强的 GAN-based 解码器，重建质量提升。
  \item \warn{架构不变}：仍然是解耦双编码器 + 共享 LLM，核心思路不变。
\end{itemize}
\end{conceptbox}

% ------- Category 6: Show-o 离散掩码扩散 -------
\subsubsection{六、Show-o 离散掩码扩散（MDLM）}

\begin{conceptbox}{Q6-1：Show-o 的 MDLM 前向/逆向过程}
\key{MDLM（Masked Diffusion Language Model）}：在\textbf{离散 token 空间}定义扩散过程。

\textbf{前向过程}：以概率 $\alpha_t$ 将 token 替换为 $[\text{MASK}]$：
\[
q(x_t | x_0) = \alpha_t \cdot \mathbf{1}[x_t = \text{MASK}] + (1-\alpha_t) \cdot \mathbf{1}[x_t = x_0]
\]

\textbf{逆向过程}：模型预测被 mask 的原始 token：
\[
p_\theta(x_0 | x_t) = \text{Transformer}(x_t)
\]
每步选择\key{置信度最高}的部分 mask token 解码（MaskGIT 策略）。

\textbf{训练损失}：仅对 $[\text{MASK}]$ 位置计算 CE loss：
\[
\mathcal{L}_{\text{MDLM}} = -\mathbb{E}_{t, x_t} \sum_{i: x_t^i = \text{MASK}} \log p_\theta(x_0^i | x_t)
\]
\end{conceptbox}

\begin{conceptbox}{Q6-2：Show-o 的 Omni-Attention 设计}
Show-o 使用\key{Omni-Attention}来统一处理文本（AR）和图像（MDLM）token：
\begin{itemize}
  \item \textbf{文本 token}：causal attention（下三角掩码）
  \item \textbf{图像 token}：full bidirectional attention（图像内部互相可见）
  \item \textbf{文本 $\leftrightarrow$ 图像跨模态}：文本 attend 到之前的文本和图像；图像 attend 到所有文本但\textbf{不 attend 到其他图像 token 序列}（避免不同图像间污染）
\end{itemize}

这一设计使单个 Transformer 可同时支持 VQA（理解）和图像生成（MDLM），无需切换模式。
\end{conceptbox}

\begin{conceptbox}{Q6-3：MaskGIT 置信度采样为什么比逐 token AR 快？}
\textbf{标准 AR}：$N$ 个图像 token 需要 $N$ 次前向（如 256 次）。

\textbf{MaskGIT 策略}：每步\textbf{并行预测所有 MASK 位置}，然后保留置信度最高的 $k$ 个（$k$ 按 cosine schedule 递增），其余重新 mask。总步数 $T \ll N$（通常 $T = 8 \sim 16$）。

\[
\text{speedup} \approx \frac{N}{T} \approx \frac{256}{12} \approx 20\times
\]

\warn{代价}：并行预测时各位置互相"不知道"对方的最终值，生成一致性略低于 AR，但实践中视觉质量差异不大。
\end{conceptbox}

% ------- Category 7: 架构设计 -------
\subsubsection{七、关键架构设计}

\begin{conceptbox}{Q7-1：MoT vs MoE 深度对比}
\begin{center}
\small
\begin{tabular}{|l|l|l|}
\hline
\textbf{维度} & \textbf{MoT（Mixture of Transformers）} & \textbf{MoE（Mixture of Experts）} \\
\hline
路由方式 & 确定性（token 类型标签） & 学习路由器（softmax top-k）\\
路由参数 & 无额外参数 & 路由网络参数 $\sim O(d)$ \\
负载均衡 & 自然均衡（按模态分配） & 需要辅助 loss（aux\_loss）\\
适用场景 & 模态明确可区分 & 任意任务混合 \\
代表模型 & BAGEL & Mixtral, DeepSeek-MoE \\
稳定性 & 高（无路由抖动） & 训练初期路由不稳定 \\
\hline
\end{tabular}
\end{center}
\end{conceptbox}

\begin{conceptbox}{Q7-2：2D RoPE 和 M-RoPE 如何为图像 patch 编码位置？}
\textbf{标准 1D RoPE}（文本）：token 位置 $p$ 对应旋转矩阵 $R_p$，维度两两配对旋转。

\textbf{2D RoPE}（图像）：将每个 patch 的\key{行位置}和\key{列位置}分别编码，特征维度一半用于行旋转，一半用于列旋转：
\[
\text{RoPE}_{2D}(q, r, c) = R_r^{(h)} \cdot R_c^{(w)} \cdot q
\]
相对距离保持\textbf{平移不变性}，斜向 patch 间内积可反映二维空间距离。

\textbf{M-RoPE}（多模态，Qwen2-VL）：为文本、图像行、图像列\textbf{分配独立的旋转分量}，同时编码时序位置和空间位置，支持视频的时空联合编码。
\end{conceptbox}

\begin{conceptbox}{Q7-3：Cross-Attention vs Self-Attention vs Q-Former 在多模态融合中的取舍}
\begin{itemize}
  \item \textbf{Q-Former（BLIP-2）}：固定数量的可学习 query token，通过 Cross-Attention 从图像特征中提取信息，再传给 LLM。\good{信息压缩高效}，\warn{但固定 query 数量限制细节}。
  \item \textbf{Self-Attention 融合（Chameleon/Transfusion）}：图像 token 直接拼接进 LLM token 序列，所有 token 在同一 Self-Attention 层交互。\good{无信息瓶颈}，\warn{序列长度增加，计算 $O(n^2)$}。
  \item \textbf{Cross-Attention 融合（Flamingo）}：LLM 层间插入 Cross-Attention 层，LLM query attend 到图像特征。\good{LLM 参数冻结可复用}，\warn{新增参数，架构改动大}。
\end{itemize}
\end{conceptbox}

% ------- Category 8: 训练技巧 -------
\subsubsection{八、训练技巧与优化策略}

\begin{conceptbox}{Q8-1：BAGEL 的多阶段训练策略（4 阶段）}
\begin{enumerate}
  \item \textbf{Stage 1——视觉编码器对齐}：冻结 LLM，只训练 SigLIP/VAE 到 LLM 的 Adapter，用大量图文对对齐特征空间。
  \item \textbf{Stage 2——统一预训练}：解冻 LLM（MoT），同时用理解数据（图文对话）和生成数据（图像重建）联合训练，建立基础能力。
  \item \textbf{Stage 3——高质量数据微调}：使用精选高分辨率图像、指令跟随数据、审美标注数据，提升生成质量和指令遵循能力。
  \item \textbf{Stage 4——RLHF 对齐}：使用 Flow-GRPO / UniGRPO，以 GenEval / 人类偏好作为奖励信号，进一步对齐文图一致性。
\end{enumerate}
\warn{关键原则}：每阶段只解冻需要更新的模块，避免灾难性遗忘。
\end{conceptbox}

\begin{conceptbox}{Q8-2：多任务损失均衡问题及解决方案}
统一模型同时优化文本 CE loss 和图像 FM/MSE loss：
\[
\mathcal{L}_{\text{total}} = \lambda_{\text{text}} \mathcal{L}_{\text{CE}} + \lambda_{\text{image}} \mathcal{L}_{\text{FM}}
\]
\warn{问题}：$\mathcal{L}_{\text{CE}} \sim O(1)$（log-prob 量级），$\mathcal{L}_{\text{FM}} \sim O(1)$（MSE 量级），但梯度尺度和优化景观差异极大。

\textbf{解决方案}：
\begin{itemize}
  \item \key{固定 $\lambda$ 比例}：根据经验或网格搜索设定（Transfusion 使用 $\lambda_{\text{image}} = 5$）
  \item \key{GradNorm}：动态调整各 loss 权重使梯度范数平衡
  \item \key{不确定性加权}（Kendall et al.）：$\mathcal{L} = \frac{1}{2\sigma_1^2}\mathcal{L}_1 + \frac{1}{2\sigma_2^2}\mathcal{L}_2 + \log\sigma_1\sigma_2$，其中 $\sigma_i$ 可学习
  \item \key{分离优化器}：文本和图像路径使用不同 learning rate（MoT 天然支持）
\end{itemize}
\end{conceptbox}

\begin{conceptbox}{Q8-3：CFG（Classifier-Free Guidance）在统一模型中如何实现？}
\key{训练阶段}：以概率 $p_{\text{drop}} = 0.1$ 将条件（文本 prompt）替换为空字符串 $\varnothing$，模型同时学习条件生成 $p(x|c)$ 和无条件生成 $p(x)$。

\key{推理阶段}：双次前向（conditional + unconditional），线性外插：
\[
\hat{v}_\theta(x_t, c) = v_\theta(x_t, \varnothing) + w \cdot \big[v_\theta(x_t, c) - v_\theta(x_t, \varnothing)\big]
\]
其中 $w > 1$ 为引导强度（guidance scale）。

\warn{在统一模型中的挑战}：文本生成不能用 CFG（会改变语义），只对图像生成 token 应用 CFG，需要在推理时\textbf{区分当前处于文本模式还是图像模式}。
\end{conceptbox}

\begin{conceptbox}{Q8-4：Flow-GRPO / UniGRPO——如何将 RLHF 引入统一生成？}
\key{核心思路}：将 GRPO（Group Relative Policy Optimization）扩展到连续扩散/Flow Matching 生成。

\textbf{GRPO for 文本}：对同一 prompt 采样 $G$ 个输出，用 reward 差值估计优势 $A_i = r_i - \bar{r}$，策略梯度更新：
\[
\nabla \mathcal{L}_{\text{GRPO}} = \mathbb{E}\left[A_i \nabla \log \pi_\theta(y_i | x)\right]
\]

\textbf{Flow-GRPO}：用 GenEval / 人类偏好作为奖励，对图像生成的\textbf{去噪轨迹}计算策略梯度，奖励信号反向传播到 Flow Matching 网络。

\good{效果}：BAGEL 经 RLHF 后 GenEval 提升约 5-8\%，文图语义一致性显著改善。

\warn{挑战}：奖励稀疏（只在最终图像计算），需要对完整去噪链做重要性采样，方差大。
\end{conceptbox}

% ------- Category 9: 综合评估 -------
\subsubsection{九、综合评估与横向对比}

\begin{conceptbox}{Q9-1：生成质量评估指标全览}
\begin{center}
\small
\begin{tabular}{|l|l|l|l|}
\hline
\textbf{指标} & \textbf{衡量维度} & \textbf{计算方式} & \textbf{局限} \\
\hline
FID & 分布相似度 & InceptionV3 特征 Fréchet 距离 & 不感知单张图质量 \\
IS & 质量 + 多样性 & $\exp(\mathbb{E}[KL(p(y|x)\|p(y))])$ & 不比较真实分布 \\
CLIP Score & 文图对齐 & $\cos(\text{CLIP}_v, \text{CLIP}_t)$ & 受 CLIP 能力限制 \\
GenEval & 细粒度指令跟随 & 6 维子任务准确率 & 覆盖场景有限 \\
DPG-Bench & 密集 prompt 生成 & 细节一致性分 & 评测成本高 \\
\hline
\end{tabular}
\end{center}

\key{面试常考}：FID 越低越好；IS 越高越好；CLIP Score 和 GenEval 越高越好。
\end{conceptbox}

\begin{conceptbox}{Q9-2：主流统一模型综合性能对比（2025 年）}
\begin{center}
\small
\begin{tabular}{|l|l|l|l|l|}
\hline
\textbf{模型} & \textbf{生成方式} & \textbf{GenEval} & \textbf{MMMU} & \textbf{亮点} \\
\hline
Chameleon & AR + VQ & 0.39 & 22.4 & 首个真统一 \\
Janus-Pro & AR + VQ（解耦） & 0.80 & 41.0 & 解耦提升生成 \\
Show-o & AR + MDLM & 0.53 & 37.0 & 离散扩散统一 \\
Transfusion & AR + Diffusion & 0.63 & — & 联合训练范式 \\
BAGEL & MoT + FM & \textbf{0.82} & \textbf{52.0+} & 目前综合最强 \\
\hline
\end{tabular}
\end{center}
\end{conceptbox}

\begin{conceptbox}{Q9-3：AR vs 扩散——速度/质量/可控性三角权衡}
\begin{center}
\small
\begin{tabular}{|l|l|l|l|}
\hline
& \textbf{纯 AR（VQ）} & \textbf{AR + FM 混合} & \textbf{纯扩散（DiT）} \\
\hline
图像质量 & 中（VQ 瓶颈） & 高（连续 latent） & 高 \\
生成速度 & 慢（$N$ 步 AR） & 中（LLM + $T$ 步 FM） & 中（$T$ 步去噪）\\
可控性（CFG） & 弱 & 强 & 强 \\
与 LLM 融合 & 天然 & 工程复杂 & 需额外设计 \\
理解能力 & 强 & 强 & 弱 \\
\hline
\end{tabular}
\end{center}

\key{面试答题策略}：如果问"选哪种架构"，先问应用场景——纯生成选扩散，理解+生成统一选 AR+FM 混合（BAGEL），快速原型选纯 AR（Janus）。
\end{conceptbox}

\begin{conceptbox}{Q9-4：开放题——如果让你设计下一代统一理解与生成模型，你会怎么做？}
\textbf{参考思路}（展示系统性思考）：

\textbf{1. 表示层}：沿用\key{解耦视觉编码器}——SigLIP for 理解，VAE for 生成，通过 Adapter 接入统一 LLM backbone。

\textbf{2. 架构层}：采用\key{MoT}（Mixture of Transformers）——共享 Attention 层保留跨模态交互，专用 FFN 减少模态干扰，无需复杂路由。

\textbf{3. 生成层}：使用\key{Flow Matching（Rectified Flow）}取代 DDPM/DDIM——ODE 轨迹更直，推理步数更少（$\sim 8$ 步），训练目标 $v = x_1 - x_0$ 简洁。

\textbf{4. 对齐层}：引入\key{Flow-GRPO} / 人类偏好 RLHF，针对文图一致性、美学质量进行强化学习微调。

\textbf{5. 扩展性}：设计支持\key{视频生成}的时序扩展——M-RoPE 编码时空位置，3D Flow Matching 处理帧间连续性。

\note{回答此类开放题时，重点展示对现有模型优劣的深刻理解，以及如何有针对性地组合现有技术解决具体痛点。}
\end{conceptbox}

\begin{warnbox}{高频易混淆概念速查}
\begin{itemize}
  \item \warn{SigLIP vs CLIP}：SigLIP 用 Sigmoid 逐对分类（不需要归一化 softmax），训练更稳定，对小批量友好；CLIP 用 softmax 对比，需大 batch size。
  \item \warn{MoT vs MoE}：MoT 按 token \textbf{类型}确定性路由（无学习参数），MoE 按 token \textbf{内容}学习路由（有路由器参数）。
  \item \warn{MDLM vs DDPM}：MDLM 在\textbf{离散 token 空间}用 MASK 操作，DDPM 在\textbf{连续像素/latent 空间}加高斯噪声。
  \item \warn{VQ-VAE vs VAE}：前者离散化（STE 有偏梯度），后者连续（重参数化无偏梯度）；前者与 LLM 天然融合，后者图像质量更高。
  \item \warn{Flow Matching vs DDIM}：FM 直接学习速度场 $v = x_1 - x_0$，DDIM 基于 DDPM 的 ODE 轨迹；FM 路径更直（OT-FM），推理更高效。
  \item \warn{GenEval vs FID}：GenEval 测指令跟随（组合生成准确率），FID 测样本分布相似度（无参考 prompt）。
\end{itemize}
\end{warnbox}

\begin{summarybox}
\textbf{统一理解与生成模型面试要点速记}：
\begin{enumerate}
  \item 三大核心矛盾（表示空间 / 注意力模式 / 训练目标）是一切设计选择的出发点
  \item VQ（离散）vs VAE（连续）：前者与 LLM 天然融合，后者图像质量更高——\key{解耦}是当前最优工程选择
  \item Transfusion 双向注意力、BAGEL MoT、Janus 解耦编码器——三种主流混合策略的核心思路要烂熟
  \item CFG 实现：训练时 dropout 条件，推理时双次前向外插
  \item 多阶段训练 + 损失均衡 + RLHF（Flow-GRPO）是当前 SOTA 模型的标配
  \item 评估：生成用 FID/GenEval，理解用 MMMU/MMBench，文图对齐用 CLIP Score
\end{enumerate}
\end{summarybox}
